%%%%%%%%%%%%%%%%%%%%%%%%%%%%%%%%%%%%%%%%%%%%%%%%%%%%%
%			HLAVIČKA								%
%%%%%%%%%%%%%%%%%%%%%%%%%%%%%%%%%%%%%%%%%%%%%%%%%%%%%
\documentclass[openany,12pt]{memoir}

\usepackage[utf8]{inputenc} 
\usepackage[czech]{babel}
\usepackage[T1]{fontenc}
\usepackage[top=1.5cm, bottom=2cm, left=2cm, right=2cm]{geometry}  % --> NASTAVENÍ OKRAJŮ
\usepackage{fancyhdr}
\usepackage{graphicx}
%\usepackage{lmodern}
\usepackage{tgschola}  % --- Silnější font
\usepackage{xwatermark}
\usepackage{xcolor}
\usepackage{changepage}
\usepackage{newunicodechar}



%%%%%% Package na zpěvník
\usepackage[full]{leadsheets}%http://mirrors.nic.cz/tex-archive/macros/latex/contrib/leadsheets/leadsheets_en.pdf   --> dokumentace	
\definesongtitletemplate{empty}{} 
\setchords{
format = \bfseries,   %tučné akordy
minor = {mi},% 
input-notation = {german},%
output-notation = {german}%
}
\definesongtitletemplate{empty}{} 

\newlength{\drop}
\newwatermark[allpages,color=red!50,angle=0,scale=2, xpos=0,ypos=0]{\includegraphics[width=5cm]{obr/pozadi2.jpg}} %--> dvojka na pozadí


%%%% Vlastní příkazy
\newcounter{Slokočet}   %Automatické číslování slok
\newcommand{\mezera}{\vspace*{0.5cm}}   %Horizontální odsazení slok
\newcommand{\stred}{5.2cm}   %%% Na zarovnání slok doprostřed, pozn. automatičtější zarovnávání na střed nejde
\newcommand{\refren}{\mezera \noindent \textbf{R:} } %refrén
\newcommand{\sloka}{\mezera \noindent \addtocounter{Slokočet}{1} \arabic{Slokočet}. } 	%sloka, která se automaticky čísluje
\newcommand{\ssloka}{\mezera \noindent}  % vlastní číslo sloky
\newcommand{\carka}{,\:}
\newcommand{\predehra}{\mezera \hspace*{-1.85cm} \textbf{Předehra:} }
\newcommand{\m}[1]{\color{white}{#1}}  %Pro akordy
\newcommand{\ap}{'} %Pro apostrof
\newcommand{\elipsa}{\kern\fontdimen3\font} %Příkaz pro lepší zacházení s výpustkami (=...); je to vpodstatě jen mezera mezi tečkama výpustky

%%%%%%%%%%%%%%%%%%%%%%%%%%%%%%%%%%%%%%%%%%%%%%%%%%%%%%
%			NÁVOD									 %
%%%%%%%%%%%%%%%%%%%%%%%%%%%%%%%%%%%%%%%%%%%%%%%%%%%%%%
%1. Věci v hlavičce IGNOROVAT
%2. Píseň psát do prostoru mezi \begin{song} a \end{song}
%3. další řádek se značí dvěma odsazeníma (= dvakrát stisknout enter)
%4. \refren vždy na začátku refrenu a \sloka na začátek sloky (automaticky se čísluje)
%  \ = alt gr + q ; [/] = alt gr f/g ; {/} = alt gr + b/n; ^ = alt gr + 3 + mezera
%Cokoliv napíšete do ^{  } se bude brát jako akord
%když se toto bude dotýkat nějakého slova (nebude mezi tím a slovem mezera)
%tak se akordy zjeví nad slovem, ale pište to před slovo
%Když se to nedotýká slova, tak akord lítá ve vzduchu a vytiskne se větší mezera
%První možnost je asi preferovanější
%5. Akordy stačí psát jen do první sloky, když se nezmění -- kytaristi to zvládnou



%\usepackage{subfiles}
%</++>

\begin{document}

\pagestyle{empty}
%\addtocontents{toc}{\protect\thispagestyle{empty}}
%\tableofcontents \thispagestyle{empty}\newpage
\newgeometry{top=1.5cm, bottom = 0cm, left = 2cm, right = 2cm}
	
\newpage

%Pro kompilaci je důležitý, 	aby ten input a \end{document} zůstali na posledních
%dvou řádcích!!!!


\begin{song}{title=\predtitle\centering Soudný den \\\large Spirituál Kvintet   \vspace*{-0.3cm}}  %% sem se napíše jméno songu a autor
\begin{centerjustified}

\begin{varwidth}[t]{0.52\textwidth}\setlength{\parindent}{0.45cm}  %Varianta č. 2 --> Dva sloupce

\sloka
^{Ami\z}Zdál se mi sen, že se nebe hroutí,

^{G\z}zdál se mi sen o poslední pouti,

^{Ami\z}zdál se mi sen, že všechno seberou ti

^{Ami\z}v~ten ^{G\,\,\,Emi\,}soudný ^{Ami}den.


\sloka
Kam běžet mám -- slunce rychle 

chladne,

měsíc na zem spadne,

moře už je na dně 

v ten soudný den.

\sloka
Stůj, nechoď dál, času už je málo, 

míň než by se zdálo,

otevři se skálo

v ten soudný den.


\sloka
Pán tě zavolá, má pro každého místo,

Pán tě zavolá, jen kdo má duši čistou.

Pán tě zavolá, sám nedokázal bys to 

v ten soudný den

\refren
Soudí, soudí, pány, slouhy,

soudí, soudí, hříšné touhy

soudí, soudí\elipsa\ldots\,aááááá


\end{varwidth}\mezisloupci\begin{varwidth}[t]{0.52\textwidth}\setlength{\parindent}{0.45cm}\vspace*{0.0cm}  % V případě varianty č.2 jde odsud text do pravé části

\sloka
Vtom se ^{Ami\z}probudíš, to byl jen sen,

^{F\z}vtom se probudíš, to byl jen sen,

^{H\z}vtom se probudíš to byl jen sen,

^{Ami\z E}jen~pouhý ^{Ami}sen.


\sloka
Zdál se mi sen o poslední pouti,

zdál se mi sen, že se nebe hroutí,

zdál se mi sen, že všechno seberou ti

v ten soudný den.


\sloka
Zdál se mi sen, já stojím na svém místě,

zdál se mi sen, mé svědomí je čisté,

zdál se mi sen, jen jedno vím jistě.

Je soudný den!



\end{varwidth}

\end{centerjustified}
\setcounter{Slokočet}{0}
\end{song}

\end{document}
